\begin{center}
{ \LARGE Constitution of University of Ottawa Computer Science Club }

%  This should be the date of the last official change, meaning the day of the
%  vote the change.  TODO: review this policy
Last edited: April 18th 2026
\end{center}

\section*{Article 1 - Name}

The club’s official names will be ``University of Ottawa Computer Science Club’’ in English
and ``Club d’informatique de l’Université d’Ottawa’’ in French.
The semi-abbreviated ``uOCSClub’’ may also be used.
No other name will be used in the advertisement or representation of the club.

\section*{Article 2 - Mission Statement}

% It should be general enough to not dictate us to perform specific activities,
% but to provide a guideline of intention behind our events and activities.

\begin{enumerate}
    % Hackathons, workshops, and other events
    \item To host and promote social and competitive hacking challenges;
    % Managing our Slack channel, partially related to our social events, notifying
    % community of important developments in IT.
    \item To foster a hacker / coder community at the University of Ottawa;
    % 
    \item To provide coding sessions and technology workshops;
    \item To organize community-driven, system administration projects; and
    \item To be a forum for teaching and sharing IT knowledge.
\end{enumerate}

\section*{Article 3 - Membership}

% This part of the constitution is important.  We need to define what it means to
% be a member of CS Club.

An individual becomes a recognized \emph{member} of the club upon official registration with the club's executives.

The executives' responsibilities and mandate in relation to the membership of the
club are:
\begin{enumerate}
    \item Ensuring that membership is open, accessible, and free to all individuals;
    \item Ensuring that a minimum of 75\% of the club's members also members of UOSU; and
    \item Requiring that all executive members of the club are members of UOSU;
\end{enumerate}
% Member rights and clauses on disability Sat 11 Oct 2025 13:58:25 EDT
\subsection*{Rights of Members}
All members of the club shall have the following rights:
\begin{enumerate}
    \item The right to vote in all elections, amendments, and referenda of the club;
    \item The right to seek election to the Executive, subject to eligibility requirements;
    \item The right to propose amendments to the Constitution, subject to Article 8;
    \item The right to access the governing documents of the club, including meeting minutes, in an accessible format upon request;
    \item The right to participate fully in club activities regardless of disability status, with reasonable accommodations provided where necessary upon request.
\end{enumerate}

\section*{Article 4 - Executives}

% Need:
%  2 people to have signing authority
%  to define our council
%  (if CSSA thing happens) official representative
%  
The club shall be governed the Executive comittee; defined here:
\begin{enumerate}
    \item Executive signing officers (max 3): \begin{enumerate}
        \item President (or Co-Presidents) 
        \item Treasurer
    \end{enumerate}
    \item Executives: \begin{enumerate}
        \item VP Logistics
        \item VP External
        \item VP CS Games
        \item VP Marketing
    \end{enumerate}
\end{enumerate}

\textit{There must at all times be a minimum of two signing officers. Other executive positions may be vacant.}

\textit{In the event that an executive position becomes vacant, its responsibilities shall be shared among the remaining members of the Executive committee until the position is filled.}

\textit{One of the signing officers will be the official club representative for the CSSA}

\section*{Article 5 - Responsibilities of the Executive}

The minimal responsibilities of the Executive are as follows:
\begin{enumerate}
    \item All Executives will:
    \begin{enumerate}
      \item Organize at least one event a semester.
    \end{enumerate}

    \item The President will:
    \begin{enumerate}
        \item Oversee the other members of the executive in fulfilling their
              responsibilities;
        \item Chair all meetings;
        \item Have signing authority for the club; and
        \item Take minutes for all meetings.
    \end{enumerate}

    \item The Treasurer will:  % Addition
    \begin{enumerate}
      \item Maintain accurate financial records of all club transactions;
      \item Manage the club’s bank account and budget allocation for events, workshops, and competitions;
      \item Handle funding applications (UOSU, sponsors, etc.) and reimbursements in accordance with university policies;
    \end{enumerate}


    \item The VP Logistics will:  % Addition
    \begin{enumerate}
      \item Coordinate logistics for events and competitions, including room bookings, equipment, and transportation;
      \item Manage inventory such as club merchandise, banners, and event materials;
      \item Ensure all events comply with university safety and booking policies;
      \item Work with VP Finance to plan event budgets and purchases.
    \end{enumerate}

    \item The VP External will:
    \begin{enumerate}
        \item Respond to club emails;
        \item Reach out to external organizations for sponsors and funding;
        \item Attend the monthly CSSA Meetings.
    \end{enumerate}

    \item The VP CS Games will:
    \begin{enumerate}
      \item Organize a yearly intro session to CS Games;
      \item Ensure coordination for CS Games training sessions;
      \item Coordinate tryouts and logistics for the yearly CS Games competition.
      \item Attend the yearly CS Games competition as team lead. 
    \end{enumerate}

    \item The VP Marketing will:
    \begin{enumerate}
        \item Post on social media (Discord, Instagram) to promote club activities;
        \item Reach out to clubs for promotional events and collaboration;
        \item Respond to personal messages on social media platforms.
    \end{enumerate}
\end{enumerate}

\section*{Article 6 - Appointed Officials}

The executives may appoint officials to aid them with their club administration tasks. These officials do not have voting rights in executive decisions. These positions may remain vacant.

The minimal responsibilities of the appointed officials are outlined here:

\begin{enumerate}
    \item The Appointed Advisors will:
    \begin{enumerate}
        \item Provide counsel and guidance to the Executive committee on club matters.
    \end{enumerate}
    \item The Appointed Director of Design will:
    \begin{enumerate}
        \item Design promotional materials for the club;
        \item Assist the VP Marketing in their duties;
        \item Lead the design of all front-facing materials for the club.
    \end{enumerate}
    \item The First Year Ambassador will:
    \begin{enumerate}
        \item Represent first-year students and communicate their concerns or suggestions to the Executive committee;
        \item Help promote club activities to incoming students and encourage participation;
        \item Assist in organizing first-year–focused events such as intro workshops or mentorship programs.
    \end{enumerate}
\end{enumerate}

\section*{Article 7 - Meetings}

The executive team is responsible for organizing one General Assembly per school year and executive meetings as required.

\begin{enumerate}
    \item General Assembly
        \begin{enumerate}
            \item It must be advertised one week before it unfolds
            \item It will take place in the end of the summer semester
            \item It is open to all registred uOCSClub members
            \item It may be in-person, online or hybrid
            \item If it is in-person, it must occur on or around the uOttawa campus
            \item Meeting minutes must be published following the meeting
        \end{enumerate}
    \item Executive meetings
        \begin{enumerate}
            \item After the appointement of the new executives, they will determine how they will be handled
            \item They do not need to be advertised
            \item Members may request information about when and where they are held
            \item They are open to all members of uOCSClub
        \end{enumerate}
\end{enumerate}

\section*{Article 8 - Elections}

% The election process is for the executive members.

\begin{enumerate}
    \item Elections for the Executive will take place in the end of the winter semester.
    \begin{enumerate}
        \item No earlier than the last week of March;
        \item No later than the last week of April.
    \end{enumerate}
    \item The length of the mandate of executive members is from their election to the end of the elections the following year.
     \item The executive positions will be open to members of the club who are also members of UOSU
    \item Any online voting system used shall be administered by a neutral platform or election officer designated by the Executive committee;
     \item All elections shall be conducted in a fair and transparent manner, ensuring equal opportunity for all members to participate;
     \item Elections shall be made accessible to all members, including the provision of reasonable accommodations for members with disabilities upon request;
     \item Where possible, elections and related materials shall be made available in both English and French.
     \item The election process shall proceed through self-nomination for a position through an online survey platform. 
     \item Members may only be nominated for one position
     \item The self-nomination phase will last for 5 days.
     \item Uncontested nominations will be elected to the position.
     \item Contested nominations will proceed as follows:
         \begin{enumerate}
             \item Contested nominees will have 48h from the end of the nomination period to campaign for their position. 
             \item Following the campaigning period, an online ballot will be cast for 24h for each contested position
             \item During the election, each club member will have one (1) vote for each executive position; 
             \item The candiate with the most votes will be elected
             \item Ties will be broken by (2/3) majority decision of the previous executives, while respecting the UOSU guidelines for fair proceedings and anti-discrimination.
             \item Following these proceedings, election results will be posted
         \end{enumerate}
     \item Candidates may at any time before being elected, withdraw their nomination 
     \item The executive committee may disqualify candidates with cause through a (2/3) majority vote so long as the cause respects UOSU guidelines for fair proceedings and anti-discrimination.
     \item Vacant positions following the election may filled through appointment by the new executive comittee while respecting the UOSU guidelines for fair proceedings and anti-discrimination.


\end{enumerate}

\section*{Article 9 - Amendments}

\begin{enumerate}
    \item Amendments to the constitution must win a two-thirds (2/3) majority vote of the members present in the meeting;

    \item At least half of the executive comittee must be present.

    \item An amendment to the constitution must be published to the members, which must be presented with a typed copy of the amendment as well as typed minutes from the meeting when the amendment was passed in order to prove that the amendment was passed.
\end{enumerate}

\section*{Article 10 - Impeachment}
\begin{enumerate}
    \item Any member of the club who commits an act negatively affecting the
        interests of the club and its members may be given notice of
        impeachment at any club meeting;
    \item The impeached individual shall have the right to defend their
        actions; and,
    \item A two-thirds (2/3) majority vote of members present will result in
        the removal of the impeached individual from the club and the loss
        of any privileges associated with the club.
\end{enumerate}

\section*{Article 11 - Executive Removal Process}

\begin{enumerate}
    \item In the event that the President of the University of Ottawa Computer Science Club engages in actions that are detrimental to the club's interests or fails to fulfill their responsibilities, the general members and executives may initiate a removal process. The process for the removal of the President shall follow the steps outlined below:
    
    \begin{enumerate}
        \item Notice of impeachment: Any member or executive can submit a written notice of impeachment to the Executive committee, detailing the specific actions or behaviours that warrant removal. The notice of impeachment should be presented during an Executive committee meeting, and a copy of the notice shall be provided to the impeached President;
        
        \item Right to defend: The impeached President shall have the right to defend themselves during a special meeting called for this purpose. This meeting shall be held within a reasonable timeframe after the notice of impeachment has been submitted. The President may present their defence and provide evidence or witnesses in their favour;
        
        \item Voting for removal: After the President has presented their defence, a vote for removal shall be conducted. A two-thirds (2/3) majority vote of the members present at the meeting, including both general members and executives, shall be required to remove the President from their position;
        
        \item Loss of privileges: If the vote for removal is successful, the impeached President shall be immediately removed from their position and lose any privileges associated with the club's presidency. The President's duties and responsibilities shall be transferred to the Vice President or another designated executive member until a new President is elected.
    \end{enumerate}
    
    
    \item If any executive member, other than the President, engages in actions that are detrimental to the club's interests or fails to fulfill their responsibilities, the general members and other executives may initiate a removal process. The process for the removal of executives shall follow a similar procedure as outlined for the removal of the President, with necessary modifications to account for the specific executive position in question.
\end{enumerate}

\section*{Article 12 - Agency Clause}

% This would need to be updated to include CSSA or other overseeing organizations.
\begin{enumerate}
    \item The University of Ottawa Computer Science Club is not an agent of the University of Ottawa and its views and actions do not represent those of the University of Ottawa.

    \item  The University of Ottawa Computer Science Club is not an agent of the University of Ottawa Students’ Union (UOSU), and its views are not representative of those of the UOSU, unless stated otherwise by the UOSU.

    \item  The University of Ottawa Computer Science Club is not an agent of the Computer Science Student Association (CSSA), and its views are not representative of those of the CSSA, unless stated otherwise by the CSSA.
\end{enumerate}

% Include discussion on club meetings

\section*{Article 13 - Equity and Accessibility}
% Added this to stay compliant with UOSU  Sat 11 Oct 2025 13:55:52 EDT
\begin{enumerate}
    \item The club affirms its commitment to equity, diversity, and inclusion in all of its operations, activities, and events;
    \item The club shall not discriminate on the basis of race, gender, sexuality, disability, religion, language, or any other protected ground under Ontario law;
    \item All meetings, elections, and activities of the club shall be made accessible to members with disabilities, with reasonable accommodations provided upon request;
    \item The club shall strive to provide bilingual communications and services in both English and French, recognizing the bilingual nature of the University of Ottawa.
\end{enumerate}
